%%%%%%%%%%%%%%%%%%%%%%%%%%%%%%%%%%%%%%%%%%%%%%%%%%%%%%%%%%%%%%%%%%%
% Standalone general university LaTeX template -- homework sheet
% One-file version: no \input files needed.
%%%%%%%%%%%%%%%%%%%%%%%%%%%%%%%%%%%%%%%%%%%%%%%%%%%%%%%%%%%%%%%%%%%

\documentclass[11pt,a4paper]{scrartcl}

%%%%%%%%%%%%%%%%%%%%%%%%%%%%%%%%%%%%%%%%%%%%%%%%%%%%%%%%%%%%%%%%%%%
% Embedded file: includes/university_metadata.tex
%%%%%%%%%%%%%%%%%%%%%%%%%%%%%%%%%%%%%%%%%%%%%%%%%%%%%%%%%%%%%%%%%%%
\newcommand{\CourseTitle}{COURSE NAME}
\newcommand{\CourseShortTitle}{COURSE SHORT NAME}
\newcommand{\DocumentKind}{DOCUMENT KIND}
\newcommand{\DocumentTitle}{DOCUMENT TITLE}
\newcommand{\DocumentSubtitle}{}
\newcommand{\AuthorName}{YOUR NAME}
\newcommand{\InstitutionName}{INSTITUTION NAME}
\newcommand{\SubmissionDate}{\today}

\newcommand{\makeuniversitytitle}{%
    \begin{titlepage}
        \centering
        \vspace*{2cm}
        {\Large \CourseTitle\par}
        \vspace{1.2cm}
        {\huge\bfseries \DocumentTitle\par}
        \ifx\DocumentSubtitle\empty\else
            \vspace{0.4cm}
            {\Large \DocumentSubtitle\par}
        \fi
        \vspace{1cm}
        {\Large \DocumentKind\par}
        \vfill
        {\large \AuthorName\par}
        \vspace{0.2cm}
        {\large \InstitutionName\par}
        \vspace{0.8cm}
        {\large \SubmissionDate\par}
    \end{titlepage}
    \pagenumbering{arabic}
}

\newcommand{\makechaptertitle}{%
    \begin{center}
        {\Large \CourseTitle\par}
        \vspace{0.5cm}
        {\huge\bfseries \DocumentTitle\par}
        \ifx\DocumentSubtitle\empty\else
            \vspace{0.25cm}
            {\Large \DocumentSubtitle\par}
        \fi
        \vspace{0.5cm}
        {\large \AuthorName \hfill \SubmissionDate\par}
    \end{center}
    \vspace{0.5cm}
}

\newcommand{\makephysiktitle}{\makeuniversitytitle}

%%%%%%%%%%%%%%%%%%%%%%%%%%%%%%%%%%%%%%%%%%%%%%%%%%%%%%%%%%%%%%%%%%%
% Document metadata
%%%%%%%%%%%%%%%%%%%%%%%%%%%%%%%%%%%%%%%%%%%%%%%%%%%%%%%%%%%%%%%%%%%
\renewcommand{\CourseTitle}{Physik IV: Kerne und Teilchen}
\renewcommand{\CourseShortTitle}{Physik IV}
\renewcommand{\DocumentKind}{Hausaufgabe, Tafelvortrag und Lernfassung}
\renewcommand{\DocumentTitle}{Aufgabenblatt 7}
\renewcommand{\DocumentSubtitle}{Radioaktivität und Kettenzerfälle}
\renewcommand{\AuthorName}{Tobias Laser}
\renewcommand{\InstitutionName}{Universität Innsbruck}
\renewcommand{\SubmissionDate}{4. Mai 2026}

%%%%%%%%%%%%%%%%%%%%%%%%%%%%%%%%%%%%%%%%%%%%%%%%%%%%%%%%%%%%%%%%%%%
% Embedded file: includes/university_packages.tex
%%%%%%%%%%%%%%%%%%%%%%%%%%%%%%%%%%%%%%%%%%%%%%%%%%%%%%%%%%%%%%%%%%%
\usepackage[margin=2.5cm]{geometry}
\usepackage[onehalfspacing]{setspace}
\usepackage{scrlayer-scrpage}

\usepackage[T1]{fontenc}
\usepackage[utf8]{inputenc}
\usepackage[ngerman]{babel}
\usepackage{lmodern}
\usepackage{microtype}
\usepackage{csquotes}

\usepackage{amsmath,amsthm,amssymb,mathtools}
\usepackage{bm}
\usepackage{icomma}
\usepackage{cancel}

\usepackage{graphicx}
\usepackage{tikz}
\usetikzlibrary{arrows.meta,calc,decorations.pathmorphing,positioning,patterns,angles,quotes}
\usepackage{pgfplots}
\pgfplotsset{compat=1.18}

\usepackage[locale=DE,space-before-unit=true,per-mode=symbol,separate-uncertainty=true]{siunitx}
\usepackage{booktabs,multirow,array,tabularx,longtable}
\usepackage{caption,subcaption}
\usepackage{placeins}
\usepackage{enumitem}
\usepackage{xparse}

\usepackage[most]{tcolorbox}
\tcbuselibrary{breakable,skins,theorems}

\usepackage[breaklinks=true,colorlinks=true,linkcolor=blue,urlcolor=blue,citecolor=blue]{hyperref}
\usepackage[nameinlink,noabbrev]{cleveref}

\clearpairofpagestyles
\ihead{\CourseShortTitle}
\ohead{\DocumentKind}
\cfoot{\pagemark}
\pagestyle{scrheadings}

\setlist[enumerate]{itemsep=0.4em,topsep=0.4em}
\setlist[itemize]{itemsep=0.25em,topsep=0.35em}
\captionsetup{font=small,labelfont=bf}

%%%%%%%%%%%%%%%%%%%%%%%%%%%%%%%%%%%%%%%%%%%%%%%%%%%%%%%%%%%%%%%%%%%
% Embedded file: includes/university_macros.tex
%%%%%%%%%%%%%%%%%%%%%%%%%%%%%%%%%%%%%%%%%%%%%%%%%%%%%%%%%%%%%%%%%%%
\newcommand{\R}{\mathbb{R}}
\newcommand{\N}{\mathbb{N}}
\newcommand{\Z}{\mathbb{Z}}
\newcommand{\Q}{\mathbb{Q}}
\newcommand{\C}{\mathbb{C}}

\newcommand{\set}[1]{\left\{ #1 \right\}}
\newcommand{\abs}[1]{\left| #1 \right|}
\newcommand{\norm}[1]{\left\lVert #1 \right\rVert}
\newcommand{\avg}[1]{\left\langle #1 \right\rangle}
\newcommand{\paren}[1]{\left( #1 \right)}
\newcommand{\bracket}[1]{\left[ #1 \right]}
\newcommand{\ol}[1]{\overline{#1}}

\newcommand{\dd}{\,\mathrm{d}}
\newcommand{\dv}[2]{\frac{\dd #1}{\dd #2}}
\newcommand{\pdv}[2]{\frac{\partial #1}{\partial #2}}
\newcommand{\pddv}[2]{\frac{\partial^2 #1}{\partial #2^2}}
\newcommand{\evalat}[2]{\left. #1 \right|_{#2}}

\newcommand{\vect}[1]{\bm{#1}}
\newcommand{\unitvect}[1]{\hat{\bm{#1}}}
\newcommand{\grad}{\vect{\nabla}}
\newcommand{\divergence}{\grad\!\cdot}
\newcommand{\curl}{\grad\!\times}

\newcommand{\half}{\frac{1}{2}}
\newcommand{\third}{\frac{1}{3}}
\newcommand{\twothirds}{\frac{2}{3}}
\newcommand{\hbarc}{\hbar c}
\newcommand{\alphaf}{\alpha}
\newcommand{\eV}{\si{eV}}
\newcommand{\MeV}{\si{MeV}}
\newcommand{\GeV}{\si{GeV}}
\newcommand{\TeV}{\si{TeV}}

\newcommand{\nuclide}[3]{\ensuremath{^{#1}_{#2}\mathrm{#3}}}
\newcommand{\isotope}[2]{\ensuremath{^{#1}\mathrm{#2}}}
\newcommand{\anti}[1]{\ensuremath{\overline{#1}}}
\newcommand{\pbar}{\anti{p}}
\newcommand{\nbar}{\anti{n}}
\newcommand{\ee}{e^-}
\newcommand{\ep}{e^+}
\newcommand{\mup}{\mu^+}
\newcommand{\mum}{\mu^-}
\newcommand{\nue}{\nu_e}
\newcommand{\numu}{\nu_\mu}
\newcommand{\nutau}{\nu_\tau}
\newcommand{\pionp}{\pi^+}
\newcommand{\pionm}{\pi^-}
\newcommand{\pionz}{\pi^0}
\newcommand{\kaonp}{K^+}
\newcommand{\kaonm}{K^-}
\newcommand{\kaonz}{K^0}

\newcommand{\diffxs}{\frac{\dd\sigma}{\dd\Omega}}
\newcommand{\totalxs}{\sigma_{\mathrm{total}}}
\newcommand{\lum}{\mathcal{L}}
\newcommand{\smand}{s}
\newcommand{\tmand}{t}
\newcommand{\umand}{u}
\newcommand{\Ecm}{E_{\mathrm{CM}}}

\newcommand{\ie}{d.\,h.}
\newcommand{\eg}{z.\,B.}
\newcommand{\wrt}{bezüglich}

\DeclareSIUnit{\barn}{b}
\DeclareSIUnit{\millibarn}{mb}
\DeclareSIUnit{\microbarn}{\micro b}
\DeclareSIUnit{\nanobarn}{nb}
\DeclareSIUnit{\fermi}{fm}
\DeclareSIUnit{\year}{a}

%%%%%%%%%%%%%%%%%%%%%%%%%%%%%%%%%%%%%%%%%%%%%%%%%%%%%%%%%%%%%%%%%%%
% Embedded file: includes/university_boxes.tex
%%%%%%%%%%%%%%%%%%%%%%%%%%%%%%%%%%%%%%%%%%%%%%%%%%%%%%%%%%%%%%%%%%%
\tcbset{
    unibox/.style={
        enhanced,
        breakable,
        sharp corners,
        boxrule=0.6pt,
        left=1.2mm,
        right=1.2mm,
        top=1mm,
        bottom=1mm,
        before skip=0.8em,
        after skip=0.8em,
        fonttitle=\bfseries
    }
}

\newtcolorbox{calculationbox}[1][]{unibox,title=Rechnung,colback=black!2,colframe=black!55,#1}
\newtcolorbox{sidecalcbox}[1][]{unibox,title=Nebenrechnung,colback=black!1,colframe=black!40,#1}
\newtcolorbox{solutionbox}[1][]{unibox,title=Lösung,colback=blue!2,colframe=blue!45!black,#1}
\newtcolorbox{answerbox}[1][]{unibox,title=Antwort,colback=green!3,colframe=green!45!black,#1}
\newtcolorbox{notebox}[1][]{unibox,title=Notiz,colback=yellow!6,colframe=yellow!45!black,#1}
\newtcolorbox{definitionbox}[1][]{unibox,title=Definition,colback=cyan!3,colframe=cyan!45!black,#1}
\newtcolorbox{warningbox}[1][]{unibox,title=Achtung,colback=red!3,colframe=red!55!black,#1}
\newtcolorbox{summarybox}[1][]{unibox,title=Zusammenfassung,colback=violet!3,colframe=violet!50!black,#1}
\newtcolorbox{formulabox}[1][]{unibox,title=Formel,colback=orange!4,colframe=orange!65!black,#1}
\newtcolorbox{taskbox}[1][]{unibox,title=Aufgabe,colback=white,colframe=black!75,#1}
\newtcolorbox{talkbox}[1][]{unibox,title=Tafelvortrag,colback=teal!3,colframe=teal!55!black,#1}

%%%%%%%%%%%%%%%%%%%%%%%%%%%%%%%%%%%%%%%%%%%%%%%%%%%%%%%%%%%%%%%%%%%
% Embedded file: includes/university_theorems.tex
%%%%%%%%%%%%%%%%%%%%%%%%%%%%%%%%%%%%%%%%%%%%%%%%%%%%%%%%%%%%%%%%%%%
\theoremstyle{definition}
\newtheorem{definition}{Definition}[section]
\newtheorem{example}{Beispiel}[section]

\theoremstyle{plain}
\newtheorem{satz}{Satz}[section]
\newtheorem{lemma}{Lemma}[section]

\theoremstyle{remark}
\newtheorem*{remark}{Bemerkung}

%%%%%%%%%%%%%%%%%%%%%%%%%%%%%%%%%%%%%%%%%%%%%%%%%%%%%%%%%%%%%%%%%%%
\begin{document}

\makeuniversitytitle
\tableofcontents
\newpage

%%%%%%%%%%%%%%%%%%%%%%%%%%%%%%%%%%%%%%%%%%%%%%%%%%%%%%%%%%%%%%%%%%%
% Homework content starts here
%%%%%%%%%%%%%%%%%%%%%%%%%%%%%%%%%%%%%%%%%%%%%%%%%%%%%%%%%%%%%%%%%%%

\section*{Hinweise zur Lernfassung}
\addcontentsline{toc}{section}{Hinweise zur Lernfassung}

Diese Fassung ist so aufgebaut, dass sie drei Zwecke erfüllt: eine saubere Abgabe, eine nachvollziehbare Tafelpräsentation und eine ausführliche Lernhilfe. Die Rechenwege sind daher bewusst etwas ausführlicher als in einer Minimalabgabe.

\begin{summarybox}[title=Zentrale Ideen des Blattes]
\begin{itemize}
    \item Radioaktive Kerne folgen dem Zerfallsgesetz $N(t)=N(0)e^{-\lambda t}$ mit $\lambda=\ln(2)/t_{1/2}$.
    \item Die Aktivität ist die Zerfallsrate: $A(t)=\lambda N(t)$.
    \item Bei einem Kettenzerfall wird die Tochtersubstanz gleichzeitig erzeugt und wieder abgebaut. Dadurch entsteht ein charakteristisches Maximum in $N_2(t)$ bzw. $A_2(t)$.
    \item Für die Plots in Aufgabe 2 ist die dimensionslose Variable $x=\lambda_1 t$ besonders praktisch.
\end{itemize}
\end{summarybox}

\newpage

\section{Aufgabe 1: Radioaktivität}

\begin{taskbox}
\begin{enumerate}[label=(\alph*)]
    \item Natürliches Uran besteht zu \SI{99.28}{\percent} aus \isotope{238}{U} mit $t_{1/2}(\isotope{238}{U})=\SI{4.468e9}{\year}$ und zu \SI{0.72}{\percent} aus \isotope{235}{U} mit $t_{1/2}(\isotope{235}{U})=\SI{7.038e8}{\year}$. Wie alt müsste die Materie des Sonnensystems sein, falls beide Isotope bei seiner Entstehung gleich häufig vorhanden waren? Wie ist das Ergebnis zu interpretieren?
    \item Ein verschlossener Behälter enthält \isotope{210}{Po} und \SI{29.62}{\micro\gram} seines Tochterelements. Die Aktivität beträgt $A=\SI{1.56e8}{\becquerel}$. Gesucht sind die Alphazerfallsgleichung und die Herstellungszeit des Poloniums. Die Halbwertszeit beträgt $t_{1/2}=\SI{138.38}{\day}$.
\end{enumerate}
\end{taskbox}

\subsection{Teil (a): Uranisotopen-Verhältnis}

\begin{definitionbox}[title=Zerfallsgesetz und Halbwertszeit]
Für ein instabiles Nuklid gilt
\[
    N(t)=N(0)e^{-\lambda t},
    \qquad
    \lambda=\frac{\ln 2}{t_{1/2}}.
\]
Die Halbwertszeit ist die Zeit, nach der nur noch die Hälfte der ursprünglichen Kerne vorhanden ist.
\end{definitionbox}

\begin{calculationbox}[title=Ansatz über den Quotienten]
Nach Aufgabenannahme waren die Anfangshäufigkeiten gleich:
\[
    N_{238}(0)=N_{235}(0)=N_0.
\]
Heute gilt deshalb
\[
    N_{238}(t)=N_0 e^{-\lambda_{238}t},
    \qquad
    N_{235}(t)=N_0 e^{-\lambda_{235}t}.
\]
Der Quotient eliminiert die unbekannte Anfangszahl $N_0$:
\[
    \frac{N_{238}(t)}{N_{235}(t)}
    =\frac{N_0 e^{-\lambda_{238}t}}{N_0 e^{-\lambda_{235}t}}
    =e^{(\lambda_{235}-\lambda_{238})t}.
\]
Der heutige Quotient ist
\[
    \frac{N_{238}(t)}{N_{235}(t)}
    =\frac{99,28}{0,72}=137,89.
\]
Damit folgt
\[
    137,89=e^{(\lambda_{235}-\lambda_{238})t},
    \qquad
    t=\frac{\ln(137,89)}{\lambda_{235}-\lambda_{238}}.
\]
\end{calculationbox}

\begin{calculationbox}[title=Einsetzen der Halbwertszeiten]
Die Zerfallskonstanten sind
\[
    \lambda_{238}
    =\frac{\ln 2}{\SI{4.468e9}{\year}}
    =\SI{1.551e-10}{\per\year},
\]
\[
    \lambda_{235}
    =\frac{\ln 2}{\SI{7.038e8}{\year}}
    =\SI{9.849e-10}{\per\year}.
\]
Also
\[
    t
    =\frac{\ln(137,89)}{\SI{9.849e-10}{\per\year}-\SI{1.551e-10}{\per\year}}
    =\SI{5.94e9}{\year}.
\]
\end{calculationbox}

\begin{answerbox}
Unter der Annahme gleicher Anfangshäufigkeit ergibt sich
\[
    \boxed{t\approx \SI{5.94e9}{\year}}.
\]
Das ist größer als das übliche Alter des Sonnensystems von ungefähr $4{,}6\cdot 10^9\,\mathrm{a}$. Die Modellannahme ``beide Uranisotope waren bei der Entstehung des Sonnensystems gleich häufig'' ist daher zu stark. Physikalisch zeigt das Ergebnis vor allem: \isotope{235}{U} zerfällt wegen seiner kürzeren Halbwertszeit deutlich schneller als \isotope{238}{U}. Eine mögliche Interpretation ist, dass die Uranisotope bereits vor der eigentlichen Bildung des Sonnensystems in Nukleosyntheseprozessen erzeugt wurden und ihre Anfangshäufigkeiten nicht exakt gleich waren.
\end{answerbox}

\begin{talkbox}[title=Tafelvortrag zu Teil (a)]
\begin{enumerate}
    \item Zerfallsgesetz und $\lambda=\ln 2/t_{1/2}$ anschreiben.
    \item Anfangsbedingung $N_{238}(0)=N_{235}(0)$ betonen.
    \item Quotienten bilden, damit $N_0$ wegfällt.
    \item Endformel $t=\ln(99{,}28/0{,}72)/(\lambda_{235}-\lambda_{238})$ herleiten.
    \item Interpretation: Ergebnis ist kein direktes exaktes Sonnensystemalter, sondern prüft die Modellannahme.
\end{enumerate}
\end{talkbox}

\subsection{Teil (b): Polonium-210 und Tochterelement}

\begin{calculationbox}[title=Zerfallsgleichung]
Beim Alphazerfall wird ein Heliumkern ausgesendet. Dadurch sinkt die Massenzahl um $4$ und die Ordnungszahl um $2$. Polonium hat $Z=84$, daher entsteht Blei mit $Z=82$:
\[
    \boxed{\nuclide{210}{84}{Po}\longrightarrow \nuclide{206}{82}{Pb}+\nuclide{4}{2}{He}}.
\]
\end{calculationbox}

\begin{calculationbox}[title=Anzahl der noch vorhandenen Poloniumkerne]
Die Aktivität ist
\[
    A(t)=\lambda N_{\mathrm{Po}}(t).
\]
Damit
\[
    N_{\mathrm{Po}}(t)=\frac{A(t)}{\lambda}.
\]
Die Halbwertszeit muss hier in Sekunden eingesetzt werden, weil die Aktivität in \si{\becquerel} = \si{\per\second} gegeben ist:
\[
    \lambda
    =\frac{\ln 2}{\SI{138.38}{\day}\cdot \SI{24}{\hour\per\day}\cdot \SI{3600}{\second\per\hour}}
    =\SI{5.798e-8}{\per\second}.
\]
Also
\[
    N_{\mathrm{Po}}(t)
    =\frac{\SI{1.56e8}{\per\second}}{\SI{5.798e-8}{\per\second}}
    =2,69\cdot 10^{15}.
\]
\end{calculationbox}

\begin{calculationbox}[title=Anzahl der Blei-Tochterkerne]
Das Tochterelement ist \isotope{206}{Pb}. Aus der Masse folgt mit $M_{\mathrm{Pb}}\approx \SI{206}{\gram\per\mole}$:
\[
    N_{\mathrm{Pb}}=\frac{m_{\mathrm{Pb}}}{M_{\mathrm{Pb}}}N_A.
\]
Einsetzen:
\[
    N_{\mathrm{Pb}}
    =\frac{\SI{29.62e-6}{\gram}}{\SI{206}{\gram\per\mole}}\cdot \SI{6.022e23}{\per\mole}
    =8,66\cdot 10^{16}.
\]
\end{calculationbox}

\begin{calculationbox}[title=Zeit seit Herstellung]
Da der Behälter verschlossen war und zu Beginn nur Polonium angenommen wird, gilt
\[
    N_0=N_{\mathrm{Po}}(t)+N_{\mathrm{Pb}}(t).
\]
Gleichzeitig gilt
\[
    N_{\mathrm{Po}}(t)=N_0 e^{-\lambda t}.
\]
Also
\[
    \frac{N_{\mathrm{Po}}(t)}{N_{\mathrm{Po}}(t)+N_{\mathrm{Pb}}(t)}=e^{-\lambda t}.
\]
Logarithmieren liefert
\[
    t=-\frac{1}{\lambda}\ln\paren{\frac{N_{\mathrm{Po}}(t)}{N_{\mathrm{Po}}(t)+N_{\mathrm{Pb}}(t)}}.
\]
Einsetzen:
\[
    t
    =-\frac{1}{\SI{5.798e-8}{\per\second}}
    \ln\paren{\frac{2,69\cdot 10^{15}}{2,69\cdot 10^{15}+8,66\cdot 10^{16}}}
    =6,04\cdot 10^7\,\mathrm{s}.
\]
Umrechnung in Tage:
\[
    t=\frac{6,04\cdot 10^7\,\mathrm{s}}{86400\,\mathrm{s/d}}=699\,\mathrm{d}.
\]
\end{calculationbox}

\begin{answerbox}
Das Polonium wurde vor ungefähr
\[
    \boxed{t\approx \SI{699}{\day}}
\]
hergestellt. Das entspricht etwa $1{,}91$ Jahren.
\end{answerbox}

\begin{talkbox}[title=Tafelvortrag zu Teil (b)]
\begin{enumerate}
    \item Zerfallsgleichung hinschreiben und erklären: $A\downarrow$ um $4$, $Z\downarrow$ um $2$.
    \item Aus der Aktivität $A=\lambda N$ die noch vorhandene Anzahl $N_{\mathrm{Po}}$ bestimmen.
    \item Aus der Bleimasse die Tochterkernzahl $N_{\mathrm{Pb}}$ bestimmen.
    \item Verschlossener Behälter: $N_0=N_{\mathrm{Po}}+N_{\mathrm{Pb}}$.
    \item Zerfallsgesetz nach $t$ auflösen.
\end{enumerate}
\end{talkbox}

\FloatBarrier
\newpage

\section{Aufgabe 2: Kettenzerfälle -- Theorie}

\begin{taskbox}
Ein Radionuklid $X_1$ mit Zerfallskonstante $\lambda_1$ zerfällt in die Tochtersubstanz $X_2$. Diese zerfällt mit Zerfallskonstante $\lambda_2$ weiter in das stabile Nuklid $X_3$:
\[
    X_1 \xrightarrow{\lambda_1} X_2 \xrightarrow{\lambda_2} X_3.
\]
\begin{enumerate}[label=(\alph*)]
    \item Differentialgleichungen für $N_1(t)$ und $N_2(t)$ formulieren und lösen, wenn bei $t=0$ nur $N_0$ Kerne $X_1$ vorhanden sind.
    \item $N_3(t)$ bestimmen.
    \item Sonderfall $\lambda_1=\lambda_2$ bestimmen.
    \item Aktivitäten in Einheiten von $\lambda_1N_0$ für $\lambda_2=\lambda_1/5$ und $\lambda_2=5\lambda_1$ im Bereich $\lambda_1t=0\dots 5$ linear und logarithmisch darstellen.
\end{enumerate}
\end{taskbox}

\subsection{Teil (a): Differentialgleichungen und Lösung für $N_1$, $N_2$}

\begin{calculationbox}[title=Differentialgleichungen]
Für $X_1$ gibt es nur Verlust durch Zerfall:
\[
    \dv{N_1}{t}=-\lambda_1N_1.
\]
Für $X_2$ gibt es Produktion aus $X_1$ und Verlust durch eigenen Zerfall:
\[
    \dv{N_2}{t}=\lambda_1N_1-\lambda_2N_2.
\]
Die Anfangsbedingungen lauten
\[
    N_1(0)=N_0,
    \qquad
    N_2(0)=0,
    \qquad
    N_3(0)=0.
\]
\end{calculationbox}

\begin{calculationbox}[title=Lösung für $N_1(t)$]
Die Gleichung
\[
    \dv{N_1}{t}=-\lambda_1N_1
\]
ist separierbar:
\[
    \frac{\dd N_1}{N_1}=-\lambda_1\dd t.
\]
Integration ergibt
\[
    \ln N_1=-\lambda_1t+C.
\]
Exponentiation und die Anfangsbedingung $N_1(0)=N_0$ liefern
\[
    \boxed{N_1(t)=N_0e^{-\lambda_1t}}.
\]
\end{calculationbox}

\begin{calculationbox}[title=Lösung für $N_2(t)$ bei $\lambda_1\neq \lambda_2$]
Einsetzen von $N_1(t)$ in die Gleichung für $N_2$ liefert
\[
    \dv{N_2}{t}+\lambda_2N_2=\lambda_1N_0e^{-\lambda_1t}.
\]
Das ist eine lineare Differentialgleichung erster Ordnung. Der integrierende Faktor ist
\[
    \mu(t)=e^{\lambda_2t}.
\]
Multiplikation damit ergibt
\[
    \dv{}{t}\paren{e^{\lambda_2t}N_2(t)}
    =\lambda_1N_0e^{(\lambda_2-\lambda_1)t}.
\]
Integration von $0$ bis $t$ und $N_2(0)=0$:
\[
    e^{\lambda_2t}N_2(t)
    =\lambda_1N_0\int_0^t e^{(\lambda_2-\lambda_1)t'}\dd t'
    =\frac{\lambda_1N_0}{\lambda_2-\lambda_1}\paren{e^{(\lambda_2-\lambda_1)t}-1}.
\]
Multiplikation mit $e^{-\lambda_2t}$ liefert
\[
    \boxed{
    N_2(t)=\frac{\lambda_1N_0}{\lambda_2-\lambda_1}
    \paren{e^{-\lambda_1t}-e^{-\lambda_2t}}
    }
    \qquad (\lambda_1\neq\lambda_2).
\]
\end{calculationbox}

\begin{answerbox}
Für $\lambda_1\neq\lambda_2$ gilt
\[
    \boxed{N_1(t)=N_0e^{-\lambda_1t}},
    \qquad
    \boxed{N_2(t)=\frac{\lambda_1N_0}{\lambda_2-\lambda_1}
    \paren{e^{-\lambda_1t}-e^{-\lambda_2t}}}.
\]
\end{answerbox}

\subsection{Teil (b): Anzahl $N_3(t)$ des stabilen Nuklids}

\begin{calculationbox}
Da $X_3$ stabil ist, verschwindet kein $X_3$ mehr. Jeder ursprünglich vorhandene Kern ist zu jedem Zeitpunkt entweder noch $X_1$, gerade $X_2$ oder bereits $X_3$:
\[
    N_1(t)+N_2(t)+N_3(t)=N_0.
\]
Daraus folgt
\[
    N_3(t)=N_0-N_1(t)-N_2(t).
\]
Einsetzen der Ergebnisse aus Teil (a):
\[
    N_3(t)
    =N_0-N_0e^{-\lambda_1t}
    -\frac{\lambda_1N_0}{\lambda_2-\lambda_1}
    \paren{e^{-\lambda_1t}-e^{-\lambda_2t}}.
\]
\end{calculationbox}

\begin{answerbox}
Für $\lambda_1\neq\lambda_2$ gilt
\[
    \boxed{
    N_3(t)
    =N_0-N_0e^{-\lambda_1t}
    -\frac{\lambda_1N_0}{\lambda_2-\lambda_1}
    \paren{e^{-\lambda_1t}-e^{-\lambda_2t}}
    }.
\]
Äquivalent und oft kompakter:
\[
    \boxed{
    N_3(t)=N_0\bracket{1+\frac{\lambda_1e^{-\lambda_2t}-\lambda_2e^{-\lambda_1t}}{\lambda_2-\lambda_1}}
    }.
\]
\end{answerbox}

\subsection{Teil (c): Sonderfall $\lambda_1=\lambda_2$}

\begin{calculationbox}[title=Grenzwert mit l'Hospital]
Ausgangspunkt ist
\[
    N_2(t)=\frac{\lambda_1N_0}{\lambda_2-\lambda_1}
    \paren{e^{-\lambda_1t}-e^{-\lambda_2t}}.
\]
Setze $\lambda_1=\lambda$ und lasse $\lambda_2\to\lambda$:
\[
    N_2(t)=\lambda N_0\frac{e^{-\lambda t}-e^{-\lambda_2t}}{\lambda_2-\lambda}.
\]
Der Bruch hat die Form $0/0$. Mit l'Hospital:
\[
    \lim_{\lambda_2\to\lambda}
    \frac{e^{-\lambda t}-e^{-\lambda_2t}}{\lambda_2-\lambda}
    =
    \lim_{\lambda_2\to\lambda}
    \frac{t e^{-\lambda_2t}}{1}
    =t e^{-\lambda t}.
\]
Damit
\[
    \boxed{N_2(t)=N_0\lambda t e^{-\lambda t}}.
\]
\end{calculationbox}

\begin{notebox}[title=Warum entsteht hier ein Faktor $t$?]
Bei gleichen Zerfallskonstanten konkurrieren Aufbau und Abbau von $X_2$ auf derselben Zeitskala. Der Faktor $t$ beschreibt den anfänglichen Aufbau aus $X_1$, während $e^{-\lambda t}$ den radioaktiven Abbau beschreibt.
\end{notebox}

\subsection{Teil (d): Aktivitäten und Diagramme}

\begin{definitionbox}[title=Aktivität und dimensionslose Variablen]
Die Aktivität eines instabilen Nuklids ist
\[
    A_i(t)=\lambda_iN_i(t).
\]
Für die Darstellung verwenden wir
\[
    x=\lambda_1t,
    \qquad
    \mu=\frac{\lambda_2}{\lambda_1}.
\]
Dann gilt für die dimensionslosen Aktivitäten
\[
    \frac{A_1}{\lambda_1N_0}=e^{-x},
\]
und
\[
    \frac{A_2}{\lambda_1N_0}
    =\frac{\lambda_2N_2}{\lambda_1N_0}
    =\frac{\mu}{\mu-1}\paren{e^{-x}-e^{-\mu x}}.
\]
Da $X_3$ stabil ist, gilt $A_3=0$.
\end{definitionbox}

\begin{answerbox}[title=Zu zeichnende Funktionen]
\[
    \boxed{\frac{A_1}{\lambda_1N_0}=e^{-x}},
    \qquad
    \boxed{\frac{A_2}{\lambda_1N_0}=\frac{\mu}{\mu-1}\paren{e^{-x}-e^{-\mu x}}},
    \qquad
    x=\lambda_1t.
\]
Die beiden geforderten Fälle sind $\mu=1/5$ und $\mu=5$.
\end{answerbox}

\begin{figure}[htbp]
\centering
\begin{tikzpicture}
\begin{axis}[
    width=0.86\textwidth,
    height=6.2cm,
    xlabel={$x=\lambda_1t$},
    ylabel={$A/(\lambda_1N_0)$},
    xmin=0,xmax=5,
    ymin=0,ymax=1.05,
    grid=both,
    legend pos=north east,
    title={Linear: $\lambda_2=\lambda_1/5$}
]
\addplot[thick,domain=0:5,samples=250]{exp(-x)};
\addlegendentry{$A_1$}
\addplot[thick,dashed,domain=0:5,samples=250]{(0.2/(0.2-1))*(exp(-x)-exp(-0.2*x))};
\addlegendentry{$A_2$}
\end{axis}
\end{tikzpicture}
\caption{Aktivitäten für $\lambda_2=\lambda_1/5$ mit linearer y-Achse.}
\end{figure}

\begin{figure}[htbp]
\centering
\begin{tikzpicture}
\begin{semilogyaxis}[
    width=0.86\textwidth,
    height=6.2cm,
    xlabel={$x=\lambda_1t$},
    ylabel={$A/(\lambda_1N_0)$},
    xmin=0,xmax=5,
    ymin=1e-3,ymax=1.2,
    grid=both,
    legend pos=south west,
    title={Logarithmisch: $\lambda_2=\lambda_1/5$}
]
\addplot[thick,domain=0:5,samples=250]{exp(-x)};
\addlegendentry{$A_1$}
\addplot[thick,dashed,domain=0.001:5,samples=250]{(0.2/(0.2-1))*(exp(-x)-exp(-0.2*x))};
\addlegendentry{$A_2$}
\end{semilogyaxis}
\end{tikzpicture}
\caption{Aktivitäten für $\lambda_2=\lambda_1/5$ mit logarithmischer y-Achse. Bei $x=0$ ist $A_2=0$ und kann auf einer logarithmischen Achse nicht dargestellt werden.}
\end{figure}

\begin{figure}[htbp]
\centering
\begin{tikzpicture}
\begin{axis}[
    width=0.86\textwidth,
    height=6.2cm,
    xlabel={$x=\lambda_1t$},
    ylabel={$A/(\lambda_1N_0)$},
    xmin=0,xmax=5,
    ymin=0,ymax=1.05,
    grid=both,
    legend pos=north east,
    title={Linear: $\lambda_2=5\lambda_1$}
]
\addplot[thick,domain=0:5,samples=250]{exp(-x)};
\addlegendentry{$A_1$}
\addplot[thick,dashed,domain=0:5,samples=250]{(5/(5-1))*(exp(-x)-exp(-5*x))};
\addlegendentry{$A_2$}
\end{axis}
\end{tikzpicture}
\caption{Aktivitäten für $\lambda_2=5\lambda_1$ mit linearer y-Achse.}
\end{figure}

\begin{figure}[htbp]
\centering
\begin{tikzpicture}
\begin{semilogyaxis}[
    width=0.86\textwidth,
    height=6.2cm,
    xlabel={$x=\lambda_1t$},
    ylabel={$A/(\lambda_1N_0)$},
    xmin=0,xmax=5,
    ymin=1e-3,ymax=1.2,
    grid=both,
    legend pos=south west,
    title={Logarithmisch: $\lambda_2=5\lambda_1$}
]
\addplot[thick,domain=0:5,samples=250]{exp(-x)};
\addlegendentry{$A_1$}
\addplot[thick,dashed,domain=0.001:5,samples=250]{(5/(5-1))*(exp(-x)-exp(-5*x))};
\addlegendentry{$A_2$}
\end{semilogyaxis}
\end{tikzpicture}
\caption{Aktivitäten für $\lambda_2=5\lambda_1$ mit logarithmischer y-Achse.}
\end{figure}

\FloatBarrier

\begin{calculationbox}[title=Lerninterpretation der Maxima]
Das Maximum von $A_2$ liegt dort, wo
\[
    \dv{}{x}\paren{\frac{A_2}{\lambda_1N_0}}=0.
\]
Aus
\[
    \dv{}{x}\bracket{\frac{\mu}{\mu-1}\paren{e^{-x}-e^{-\mu x}}}=0
\]
folgt
\[
    -e^{-x}+\mu e^{-\mu x}=0,
    \qquad
    e^{-x}=\mu e^{-\mu x}.
\]
Daher
\[
    x_{\max}=\frac{\ln\mu}{\mu-1}.
\]
Für $\mu=1/5$ ergibt sich
\[
    x_{\max}=\frac{\ln(0{,}2)}{-0{,}8}=2{,}01,
\]
für $\mu=5$ dagegen
\[
    x_{\max}=\frac{\ln(5)}{4}=0{,}40.
\]
\end{calculationbox}

\begin{summarybox}[title=Physikalische Deutung von Aufgabe 2(d)]
\begin{itemize}
    \item $A_1$ fällt immer exponentiell, weil $X_1$ nur zerfällt.
    \item $A_2$ startet bei $0$, weil zu Beginn kein $X_2$ vorhanden ist.
    \item Für $\lambda_2=\lambda_1/5$ zerfällt $X_2$ langsam. Deshalb baut sich die Tochteraktivität langsam auf, erreicht ihr Maximum spät und bleibt lange relevant.
    \item Für $\lambda_2=5\lambda_1$ zerfällt $X_2$ schnell. Deshalb erreicht $A_2$ sehr früh ein Maximum und folgt danach näherungsweise dem Nachschub aus $X_1$.
    \item Auf der logarithmischen Achse sieht man besonders gut, welche Exponentialfunktion im Spätzeitverhalten dominiert.
\end{itemize}
\end{summarybox}

\begin{talkbox}[title=Tafelvortrag zu Aufgabe 2]
\begin{enumerate}
    \item Zerfallskette $X_1\to X_2\to X_3$ zeichnen.
    \item DGLs aus Produktion minus Verlust aufstellen.
    \item $N_1$ zuerst lösen; danach $N_1$ in die DGL für $N_2$ einsetzen.
    \item Für $N_2$ integrierenden Faktor zeigen, aber nicht zu lange bei Algebra bleiben.
    \item $N_3$ über Kernerhaltung bestimmen: $N_3=N_0-N_1-N_2$.
    \item Sonderfall $\lambda_1=\lambda_2$ als Grenzwert erklären: l'Hospital führt zu $N_2=N_0\lambda t e^{-\lambda t}$.
    \item Für die Plots zuerst dimensionslos machen: $x=\lambda_1t$, $\mu=\lambda_2/\lambda_1$.
\end{enumerate}
\end{talkbox}

\newpage

\section{Prüfungs- und Lernzusammenfassung}

\begin{formulabox}[title=Wichtige Formeln]
\[
    N(t)=N(0)e^{-\lambda t},
    \qquad
    \lambda=\frac{\ln 2}{t_{1/2}},
    \qquad
    A(t)=\lambda N(t).
\]
\[
    N_1(t)=N_0e^{-\lambda_1t}.
\]
\[
    N_2(t)=\frac{\lambda_1N_0}{\lambda_2-\lambda_1}
    \paren{e^{-\lambda_1t}-e^{-\lambda_2t} }
    \qquad (\lambda_1\neq\lambda_2).
\]
\[
    N_2(t)=N_0\lambda t e^{-\lambda t}
    \qquad (\lambda_1=\lambda_2=\lambda).
\]
\[
    \frac{A_1}{\lambda_1N_0}=e^{-x},
    \qquad
    \frac{A_2}{\lambda_1N_0}=\frac{\mu}{\mu-1}\paren{e^{-x}-e^{-\mu x}},
    \qquad
    x=\lambda_1t,
    \quad
    \mu=\frac{\lambda_2}{\lambda_1}.
\]
\end{formulabox}

\begin{warningbox}[title=Typische Fehler]
\begin{itemize}
    \item Aktivität in \si{\becquerel} bedeutet \si{\per\second}; deshalb muss die Halbwertszeit für $A=\lambda N$ in Sekunden eingesetzt werden.
    \item Bei Quotienten von Isotopenhäufigkeiten kürzt sich die unbekannte Anfangsmenge nur, wenn die Anfangsbedingung korrekt verwendet wird.
    \item Bei $N_2(t)$ muss das Vorzeichen stimmen. Die Formel ist positiv, obwohl der Nenner je nach Verhältnis $\lambda_2/\lambda_1$ positiv oder negativ sein kann.
    \item Eine logarithmische Achse kann den Wert $0$ nicht darstellen. Deshalb startet die logarithmische Kurve von $A_2$ praktisch bei einem kleinen positiven $x$.
\end{itemize}
\end{warningbox}

\begin{summarybox}[title=Was man für die Prüfung können sollte]
Man sollte aus dem Zerfallsgesetz mit Halbwertszeiten rechnen, Aktivitäten in Kernzahlen umrechnen, aus Tochterprodukten auf eine vergangene Herstellungszeit schließen und lineare Differentialgleichungen für einfache Zerfallsketten lösen können. Besonders wichtig ist das Verständnis, dass eine Tochtersubstanz gleichzeitig produziert und abgebaut wird. Genau deshalb steigt ihre Aktivität zuerst an und fällt später wieder ab.
\end{summarybox}

\end{document}
